\section{Terminology, Units of Analysis, and Research Questions}
\label{sec:terminology-questions}

\subsection{Attachment and processing context}

An \emph{attachment} is an application-level object submitted for transmission
or obtained after authenticated decryption. The attachment is represented by a
processing tuple
\begin{equation}
  u = (x,n,m,p,r,d),
  \label{eq:processing-unit}
\end{equation}
where $x\in\{0,\ldots,255\}^{*}$ is the source octet string, $n$ is the caller-supplied
display name, $m$ is the caller-supplied media type, $p$ is a versioned local
policy, $r$ is a request context, and $d\in\{\mathsf{out},\mathsf{in}\}$
identifies the outgoing or incoming boundary. The request context contains
only information explicitly admitted by policy, including tenant selection,
cancellation state, and trace identifiers. Neither $n$ nor $m$ is treated as a
trusted assertion about the grammar of $x$.

The \emph{source object} is the complete octet string $x$. A
\emph{candidate output} contains bytes, a generated display name, a declared
output type, the handler profile, the performed reconstruction level, and an
evidence record. It exists before independent validation. A \emph{deliverable
output} is the byte component of a candidate for which the final acceptance
predicate is true and whose public interface is available to the caller. This
distinction is central.
Creation of a temporary file, allocation of an output buffer, or successful
return from an external codec does not make an object deliverable.

A \emph{carrier} is an attachment considered with respect to a named delivery
property rather than inferred post-execution behavior. The property may be a
misleading name, an executable signature, a second grammar, an embedded active
object, excluded metadata, a malformed structural field, or a resource demand.
Calling a sample a Trojan carrier means that the sample represents a
Trojan-delivery scenario in the corpus. It does not mean that \AFD\ has inferred
Trojan behavior from the file.

\subsection{Format profiles and canonical grammars}

A \emph{format family} is a broad technical grouping such as ISO base media file
format, Ogg, or static raster image. A \emph{format profile} is narrower. It
fixes the admissible signature, structural grammar, component locations,
ordering constraints, codec set where relevant, resource bounds, output media
type, and filename suffixes. For example, a constrained WebM profile and a
general Matroska family are not interchangeable. Static PNG and APNG likewise
require different profiles even though both begin with the PNG signature.
Static and animated WebP similarly share RIFF and WebP identifiers but require
different chunk grammars and reconstruction claims.

For a format profile $k$ and policy $p$, let
\begin{equation}
  L_{p,k}\subseteq\{0,\ldots,255\}^{*}
  \label{eq:term-canonical-language}
\end{equation}
denote the accepted output language. Membership is established by a complete
parser that begins at offset zero, validates all profile constraints, and
reports a terminal offset equal to the output length. The term
\emph{canonical} in this manuscript therefore means \emph{closed under the
declared profile and policy}. It does not necessarily imply that one
mathematical byte string is the unique representation of every media object.
Where an encoder is deterministic and a unique representation is in fact
enforced, the stronger statement is named separately.

The profile is location-aware. A component that is permitted in one parent
element or phase of a grammar is not automatically permitted elsewhere. This
requirement matters for chunked and nested formats. An allowlist of FourCC,
chunk, atom, or element identifiers without parent and cardinality constraints
does not define $L_{p,k}$. The parser must also validate sizes,
offsets, ordering, cross-references, multiplicities, and exact termination.

\begin{definition}[Complete profile acceptance]
For registry parser $\mathsf{Parse}_{v,k}$ and output $y$, complete profile acceptance holds when
\begin{equation}
  \mathsf{Parse}_{v,k}(y)\Downarrow\rho_k(y),\qquad
  \rho_k.o=0,\qquad \rho_k.t=|y|,\qquad
  y\in L_{p,k},
  \label{eq:complete-profile}
\end{equation}
where $\rho_k$ is the parser report. A successful prefix parse with
$\rho_k.t<|y|$ is not acceptance.
\end{definition}

\subsection{Classification, routing, and acceptance}

The manuscript uses \emph{classification} only for evidence extraction and
routing. Three routing sources are distinguished. $E_p(n)$ is profile evidence
from the normalized final suffix. $M_p(m)$ is evidence from the supplied media
type. $B_p(x)$ is the finite candidate set returned by bounded signature and
structural-header probes. A separate complete-parser registry observes the
reconstructed output from offset zero through its exact end. It is an acceptance
observer, not a fourth routing source. No source can erase a recorded
disagreement. Unresolved disagreement prevents delivery in the strict profile.

\emph{Routing} selects a candidate handler from bounded evidence. Routing may
occur before a complete parse because the complete parser is profile-specific.
\emph{Acceptance} occurs only after reconstruction and independent output
validation. No extension, media type, magic value, decoder success, or handler
declaration can authorize delivery alone. This separation prevents a familiar
category error in file-security systems: using the mechanism that selected a
parser as evidence that the parser consumed the same object later exposed to
the application.

\subsection{Reconstruction terminology}

\begin{definition}[Structural reconstruction]
Structural reconstruction parses a container, checks its declared structural
relations, emits only allowed components, rewrites dependent lengths or
offsets, excludes policy-forbidden metadata, and accounts for the complete
source object. Encoded media payload may remain byte-identical.
\end{definition}

\begin{definition}[Semantic reconstruction]
Semantic reconstruction decodes media into a semantic intermediate
representation and encodes a new output from that representation. Source
container and codec payload bytes are not copied into the output except for
values reconstructed explicitly from validated semantic state.
\end{definition}

\begin{definition}[Semantic-or-block requirement]
Under a semantic-or-block requirement, the absence, failure, timeout, or
rejection of the semantic path produces no deliverable output. Structural
reconstruction is not an admissible fallback.
\end{definition}

The reconstruction levels form the order
\begin{equation}
  \Unprocessed \prec \Structural \prec \Semantic.
  \label{eq:term-reconstruction-order}
\end{equation}
The level is a report of work actually completed, not a property inferred from
the selected handler. A process that intended to transcode but returned the
source object has performed \Unprocessed. A container remultiplexer that copies
validated compressed frames has performed \Structural. A pixel or sample
decode followed by new encoding has performed \Semantic.

The manuscript avoids using \emph{sanitization} as an unqualified synonym for
all three outcomes. The term is used only with a named property. Structural
reconstruction may sanitize a location-metadata field. Semantic reconstruction
may neutralize a container-level appended payload. Neither statement proves
that all information encoded in retained media semantics has been removed.

\subsection{Verdicts and publication}

The externally visible verdict set is
\begin{equation}
  \mathcal{V} =
  \{\Clean,\Suspicious,\Blocked,\NotProcessed\}.
  \label{eq:verdict-space}
\end{equation}
A \Clean\ verdict indicates that the configured acceptance predicate succeeded
for the reconstructed object. \Suspicious\ records evidence that requires
policy action but is not represented as accepted delivery.
\Blocked\ records a policy rejection. \NotProcessed\ records that the required
method was not completed. This verdict tag is distinct from the performed
reconstruction level \Unprocessed. Under the evaluated release profile, only \Clean\
can carry a deliverable output.

The distinction between verdict and error is also explicit. A typed terminal
error can describe malformed input, exhausted resources, cancellation, failed
external processing, or internal invariant failure. A terminal error and a
non-clean verdict share the same publication consequence: no deliverable bytes.
This equivalence concerns the output boundary, not diagnostic meaning. Audit
records preserve the difference without exposing attacker-controlled
diagnostics as trusted application text.

\subsection{Units of measurement}

The primary security unit is one processing tuple $u$ under one pinned policy
and software build. Paired comparisons use the same source object under
predeclared policy variants. Corpus-level rates are computed over explicitly
named strata and do not mix benign usability with malicious-carrier
neutralization.

Ground truth describes how a record was constructed and which carrier property
it contains. The policy oracle is separate. It states the expected action under
one pinned policy and cites the rule that produces that action. A behavioral
malware label cannot substitute for either field.

Primary false blocking is defined only over independently usable S0 and S1
records whose policy oracle is expected-clean. Let this set be $\mathcal{B}_{01}$:
\begin{equation}
  \mathit{FBR} =
  \frac{|\{u\in\mathcal{B}_{01}:\operatorname{verdict}(u)\neq\Clean\}|}
       {|\mathcal{B}_{01}|}.
  \label{eq:term-false-blocking}
\end{equation}
Other expected-clean strata are summarized as expected-action agreement and do
not enlarge the primary benign denominator. Let $\mathcal{B}_{\mathrm{block}}$
contain records whose policy oracle requires withholding. False acceptance is
\begin{equation}
  \mathit{FAR} =
  \frac{|\{u\in\mathcal{B}_{\mathrm{block}}:
    \operatorname{verdict}(u)=\Clean\}|}
       {|\mathcal{B}_{\mathrm{block}}|}.
  \label{eq:term-false-acceptance}
\end{equation}
Carrier-property retention among expected-clean reconstructions is reported as
a different endpoint. It does not alter the expected-action denominator. A
cell with no eligible records is represented by a null value meaning ``not
estimable,'' never by zero.

Performance is measured end to end from API entry to terminal result and by
named internal stages. The unit includes classification, reconstruction,
validation, signature policy, and publication bookkeeping. Isolated decoder
time is reported separately. Latency summaries include median and upper
quantiles with confidence intervals from the pinned server run. Throughput is
not inferred as the reciprocal of an isolated median when concurrency,
allocation, and process scheduling affect the system.

\subsection{Research questions}

The study is governed by six research questions. They are stated before the
server evaluation so that corpus construction and measurement cannot be
adjusted to a preferred result.

\begingroup
\small
\setlength{\tabcolsep}{4pt}
\begin{longtable}{L{0.08\textwidth}L{0.33\textwidth}L{0.50\textwidth}}
\caption{Research questions, evidence, and decision criteria.}
\label{tab:research-questions}\\
\toprule
ID & Question & Required evidence and interpretation \\
\midrule
\endfirsthead
\toprule
ID & Question & Required evidence and interpretation \\
\midrule
\endhead
\midrule
\multicolumn{3}{r}{Continued on next page}\\
\endfoot
\bottomrule
\endlastfoot
RQ1 &
Can three routing sources and a separate complete-output observer be composed
without allowing a weak source to authorize delivery? &
Classifier-disagreement corpus, deterministic decision records, routing traces,
and proof argument for the acceptance predicate. Any accepted unresolved
disagreement refutes the release claim. \\
\addlinespace
RQ2 &
Does the reconstruction-level order prevent structural fallback when a policy
requires semantic reconstruction? &
Paired policy cases for every supported family, unavailable-decoder cases,
timeouts, cancellation, and handler-level reports. One deliverable structural
output under semantic-or-block refutes the claim. \\
\addlinespace
RQ3 &
Do accepted outputs satisfy complete consumption, profile closure, type
convergence, and declared metadata exclusions? &
Independent output parsing, differential validators, property tests,
format-specific mutations, and exact terminal offsets. Results are reported per
profile rather than as one aggregate parser-success rate. \\
\addlinespace
RQ4 &
How does the method affect malware-carrier scenarios without claiming
behavioral malware classification? &
Safe carrier simulations mapped to explicit properties, including masquerading,
prefix and appended grammars, embedded active objects, malformed structures,
and excluded metadata. Conclusions remain property-specific. \\
\addlinespace
RQ5 &
Does every failure path preserve clean-only publication and declared resource
bounds? &
Fault injection, cancellation, child-process failure, panic containment,
oversized output, diagnostic flooding, temporary-file inspection, and
interface-level invariant tests. No non-clean terminal state may expose a
deliverable artifact. \\
\addlinespace
RQ6 &
What latency, throughput, memory, expansion, and fidelity costs does the method
incur under a reproducible server configuration? &
Pinned hardware and software manifest, warm-up and repeated observations,
stratification by format and size, paired source/output fidelity measures, raw
data retention, and uncertainty intervals. No development-machine values enter
the result. \\
\end{longtable}
\endgroup

\subsection{Falsification conditions}

The principal claims are intentionally vulnerable to counterexample. The
profile-convergence claim fails if any accepted output retains an unresolved
conflict among its output name, declared media type, byte recognition, selected
handler, and complete parser. The complete-consumption claim fails if an
accepted parser report terminates before the output length. Canonical component
closure fails if an accepted output contains a component outside the
location-aware allowlist or violates a declared cardinality or ordering rule.
Metadata exclusion fails on the first accepted forbidden metadata component.

Reconstruction-strength enforcement fails if a weaker performed level satisfies
a stronger requirement. Clean-only publication fails if any non-clean verdict
or terminal error exposes an output through Rust, the filesystem, a serialized
interface, or an application wrapper. Resource-boundedness fails if the
declared limit can be exceeded before cancellation or rejection on the
platform for which the bound is claimed. Reproducibility fails if the reported
tables cannot be regenerated from the pinned source, dependency lock, corpus
manifest, policy document, and raw observations.

The malware-carrier interpretation has a different falsification boundary.
Failure to remove or reject the named carrier property is evidence against the
corresponding scenario claim. Failure to identify a malware family is not,
because family classification is not a system output. Likewise, successful
reconstruction of a media object is not evidence that the visual or acoustic
meaning is benign.

\subsection{Statement of novelty}

The scientific novelty lies in the composed decision method and its evidence
discipline, not in a new image codec, malware signature, or transport protocol.
The method joins symmetric endpoint placement around an end-to-end encrypted
channel with a ranked but non-erasing evidence model, an explicit order of
reconstruction strength, independently checked complete profile acceptance,
and transactional clean-only publication. The resulting construction turns
attachment delivery into a conjunction of separately falsifiable conditions.

The contribution differs from a conventional scanner pipeline in two ways.
First, classifier output selects work but cannot authorize publication.
Second, the representation authorized for use is the independently validated
reconstruction, not the source object on which a detector happened to run. It
differs from a conventional content-reconstruction service by preserving the
blind-relay model and by refusing to treat structural remultiplexing as
equivalent to semantic reconstruction. It differs from a parser-only method by
including policy order, resource enforcement, failure semantics, and
cross-interface publication as part of the security construction.

This novelty is methodological and systems-oriented. Its value must be judged
by the precision of the model, the traceability of implementation to the
acceptance predicate, and the pre-specified evidence contract that can refute each
claim. The extended manuscript is structured accordingly.
